\begin{papercard}{lin2025}
\cf{Stages} \stagetag{D} \stagetag{W} \stagetag{I} (also \stagetag{A}). A cross-sectional review of 115 Chinese mental-health apps on three national stores. A 61-indicator, standard-derived policy compliance scale and MARS were each applied by two calibrated raters, with very high ICCs.

\cf{Research questions} (1) How far do the privacy policies match China's PI Specification (GB/T 35273-2020)? (2) Which compliance dimensions correlate with app quality? \pc{lin2025}{3}{To what extent do the privacy policies of mental health care}.

\cf{Search} Following PRISMA 2020 \pc{lin2025}{3}{This study followed the PRISMA (Preferred}. Stores were chosen to reflect real national distribution: Google Play is unavailable in mainland China, so the Apple App Store, Huawei AppGallery and Tencent Appstore were used \pc{lin2025}{3}{the Google Play Store remains inaccessible}. Nine keywords from a preliminary test search, entered in Chinese \emph{anonymously, without logging into any account}, with all results collected \pc{lin2025}{3}{we entered these keywords anonymously in the}. For apps on both platforms with identical design and content, the Android version was chosen. Timing was based on evidence (at least 15 minutes per evaluator) and on a pilot (22.4 minutes per app on average), so 15--30 minutes was allowed per app \pc{lin2025}{3}{Preliminary data show}.

\cf{Definitions and eligibility} An explicit functional definition of mental-health care apps (structured interventions, symptom management or therapeutic assistance) and of personal information, taken from the standard \pc{lin2025}{3}{we defined mental health care apps as digital tools}. Inclusion (Textbox 1): general population; professional interventions \emph{or} disease specificity \emph{or} dynamic adaptability; collects PI at least once; free download; simplified Chinese. Exclusion: functional irrelevance; pseudoscience; generic tools only (mood diaries, meditation, ambient audio); general health management (sleep monitoring, cycle tracking); professional or indirect users; external device required; malfunction; \emph{no accessible or complete policy}; special credentials; mandatory payment or trial-only \pc{lin2025}{4}{Apps offering only generic tools}.

\cf{Material capture} Two independent screeners used names, screenshots and descriptions. Included apps were downloaded, privacy policies \emph{archived as text files}, and download and update time stamps recorded \pc{lin2025}{4}{archiving privacy policies as text files and recording download}. Flow: 4{,}583 hits, 2{,}068 unique per store, 1{,}759 after cross-store de-duplication, 273 downloaded. Exclusions: 60 technical, 9 no PI collection, 31 inaccessible policy, 13 institutional codes; 115 included \pc{lin2025}{5}{The initial search yielded}.

\cf{Instrument} A compliance scale by Ni et al.\ derived from the PI Specification: 6 level-1 indicators (PI life-cycle stages), 22 level-2 and 61 level-3 indicators, with descriptions, example sentences and the corresponding article for each (Multimedia Appendix 1). Level 3 is scored 1/0, level 2 and level 1 are averages, and the overall score is a percentage \pc{lin2025}{4}{Finally, 6 level 1 indicators, 22 level 2 indicators, and 61}.

\cf{Rater training and reliability} Systematic training by professional IT developers on technical and legal terms and on scoring. Calibration on 10 apps gave ICC 0.996, with discrepancies resolved with technical experts to set unified criteria. The main assessment then had ICC 0.988 \pc{lin2025}{5}{Before conducting the privacy policy compliance assessments}. MARS was rated by trained raters with a qualified third-person arbitrator, using a two-way random, absolute-agreement ICC \pc{lin2025}{5}{To ensure methodological rigor, ICCs based on a 2-way random}.

\cf{Key findings} All apps linked a policy in the store. 90.4\,\% showed a policy reminder at log-in, but only 73.9\,\% used a separate pop-up or click-to-confirm, and only 19.1\,\% placed the policy on the personal home page \pc{lin2025}{7}{Regarding privacy policies, it is notable that all apps provided}. Mean compliance was 60.8\,\%; lowest for information destruction (49.6\,\%). Only 2 apps notified third parties to delete data \pc{lin2025}{9}{of the apps indicated that they notified third}. Only 20.9\,\% accepted legal responsibility for security incidents \pc{lin2025}{9}{explicitly accepted legal responsibility for}. MARS correlated with compliance ($r=0.354$).

\cf{Limitations stated} Only the PI Specification (not the 2021 PIPL); legal adherence rather than effectiveness of safeguards; 15--30 minutes per app; free apps only \pc{lin2025}{14}{our evaluation metrics were exclusively}.

\cf{Use for this thesis} Strong \emph{procedural} model: anonymous search without login, storefronts chosen to match real distribution, a pilot-timed evaluation budget, policies archived as text with time stamps, calibration on 10 apps before independent coding, and reliability reported before and after. Its distinction between a policy reminder at log-in and active confirmation is directly reusable as an observable walkthrough item. Its exclusion of sleep monitoring and cycle tracking is the \emph{opposite} of the thesis scope, so cite it only as a contrasting rule. The instrument is tied to Chinese law and is not a GDPR checklist.
\end{papercard}
